Since @@M0@@ by 1.7, @@M1@@ is @@M2@@-regular by 6.2; thus @@M3@@ for @@M4@@ and @@M5@@ by 1.4 (i). The exact sequence @@M6@@ then shows that @@M7@@ is @@M8@@-regular; hence (i).

Assertion (i) implies that @@M9@@; hence (ii). By 2.10 and 1.9,
@@M10@@
thus, @@M11@@. By induction on @@M12@@, we may suppose that @@M13@@ is a @@M14@@-sheaf; hence (iii) by 1.6.2.

\sgaStableTarget{sga6.e13.lem.6-5}\textbf{Lemma 6.5.} \emph{Under the conditions 6.1, suppose that @@M15@@ has no associated prime cycle of dimension @@M16@@.}
\begin{enumerate}[label=(\roman*)]
\item \emph{If @@M17@@, then @@M18@@.}
\item \emph{@@M19@@ for @@M20@@.}
\item \emph{Let an exact sequence @@M21@@ and an integer @@M22@@ satisfy @@M23@@; then @@M24@@ and @@M25@@.}
\end{enumerate}
Indeed, let @@M26@@, @@M27@@; then @@M28@@ defines a nonzero subsheaf @@M29@@ of @@M30@@, so @@M31@@ is @@M32@@. Take a commutative diagram with exact rows and columns
@@M33@@
